\chapter{Preliminaries}
\label{ch:preliminaries}

The opening paragraph of a chapter should briefly outline the chapter structure 
by naming each section and stating what it covers.

In the preliminaries you can name the studied formalism (e.g., Model Checking,
probabilistic programs, parameter synthesis) and cite
standard references for further reading (e.g., \cite{baier2008principles}). 
Do not add a citation to every definition.

%-----------------------------------------------------------------------
\section{Basic Definitions}\label{sec:basic-definitions}
%-----------------------------------------------------------------------

Introduce the formal objects of the theory bottom-up: start with the most
elementary notion (e.g., transition systems, markov chains) and build towards the
property of interest (e.g., LTL properties, complexity).

Each definition in this section should be accompanied by:
\begin{enumerate}
    \item A sentence stating what the definition formalizes and why it is needed.
    \item The formal statement in a Definition box.
    \item A brief explanation of each component of the definition.
    \item At least one example.
\end{enumerate}

Cover the following in order:
\begin{enumerate}
    \item \textbf{Syntactic Objects.} Define the elementary building blocks
          of the formalism (e.g., lattices, markov models, program graphs), relationships between
          the individual objects and relevant formal properties of interest.
    \item \textbf{Derived Notions.} Define derived concepts such as strategies,
          reachability, bisimulations, etc.
    \item \textbf{The Property of Interest.} Give the formal definition of
          the property the thesis analyzes (e.g., termination, equivalence,
          strategy synthesis, almost-sure termination).
\end{enumerate}

%-----------------------------------------------------------------------
\section{Known Results from the Literature}\label{sec:known-results}
%-----------------------------------------------------------------------

Present the lemmas and theorems from the literature that are used in
\cref{ch:mainmatter}. 
You do not have to prove these results unless the proof is insightful for the 
rest of the thesis. 
Always cite the original source.
Each result in this section should be presented as follows:
\begin{enumerate}
    \item A sentence explaining why the result is important and how it will
     be used in the thesis.
     \item The formal statement in a Lemma or Theorem box, with a citation in
     the title.
     \item A proof indication: either a short proof, ``Proof: see 
     \cite{Doe2024Example}'', or a forward reference to the appendix.
     \item (Optional) A brief remark on the scope or limitations of the result
     if that is relevant for understanding its use in \Cref{ch:mainmatter}.
\end{enumerate}

This section collects known results from the literature that are used in the
rest of the thesis.
Organize them as follows:
\begin{enumerate}
    \item \textbf{Auxiliary lemmas.} State small technical results about the
          basic objects from \Cref{sec:basic-definitions} (e.g., equivalence relations,
          closure conditions, or monotonicity).
          These are typically needed in the proofs of \Cref{ch:mainmatter}.

    \item \textbf{Characterization theorem(s).} State the main theoretical
          results from the literature about the property of interest.
          These serve as a reference point for understanding what is already
          known.

    \item \textbf{Existing automated technique(s).} If the thesis extends an
          existing automated method (e.g., loop-invariant computation, value iteration,
          parameter synthesis), state its correctness theorem here.
\end{enumerate}

% --- Example entries (replace with results relevant to your topic) ---

\begin{lemma}[Name of Auxiliary Lemma, \cite{Doe2024Example}]
    \label{auxlemma}
    A Lemma is a small result or auxiliary result, which is only interesting because 
    we can use it to prove a more significant main result.
\end{lemma}

\begin{proof}
    The complete proof can be found in \cite{Doe2024Example}.
\end{proof}

\begin{theorem}[Name of the Key Lemma, \cite{Doe2024Example}]
    \label{keylemma}
    A Theorem is one of the main results. Everything else is defined and proven such
    that this results can be proven.
\end{theorem}

\begin{proof}
    The complete proof can be found in \cite{Doe2024Example}.
\end{proof}

\begin{example}
    Apply the theorem above to a concrete example.
    Demonstrate both a case where the result applies successfully and,
    if possible, a case where it fails due to the restriction identified above.
    The failing case motivates the extension in \cref{ch:mainmatter}.
\end{example}

\begin{corollary}[Name of a direct consequence, \cite{Doe2024Example}]
    \label{keycor}
    A Corollary is a (easy) consequence of some other (main) result. 
    Still, it has to be proven, even if the proof is rather simple!
\end{corollary}

\begin{proof}
    The complete proof can be found in \cite{Doe2024Example}.
\end{proof}

% Add further Lemma and Theorem boxes as needed.

You can use \textcolor{ForestGreen}{colors} to highlight differences in boxes 
or even difficult parts in proofs. 
However, do not make the colors too distracting and everything has to be
readable in black and white as well.